\section{Preliminaries}
\label{sec:preliminaries}

A vision--language model receives an image $I$ and question $q$ and generates an answer token by token. We write $x=(I,q)$ for this input and $a=(a_1,\ldots,a_T)$ for a fixed target answer of $T$ tokens. At layer $\ell$ and visual or text position $i$, the residual-stream state $\vec{h}_i^\ell$ is the model's complete internal vector at that location; the model uses the image, question, and preceding answer tokens to predict each next token. We call this full vector a \emph{dense} state \citep{elhage2021transformerCircuits}.

Sparse transcoders offer a second, separately fitted view of that computation. A per-layer transcoder (PLT) approximates one feed-forward update with a small set of active features \citep{dunefsky2024transcoders}; a cross-layer transcoder (CLT) can describe contributions to later layers \citep{ameisen2025circuittracing}. Neither is a component the base model uses to produce its answer. Our question is therefore twofold: does masking a human-annotated answer-bearing region lower the target-answer score more than masking comparable non-evidence content, and which specified internal object reproduces the score change caused by the evidence mask? A dense state, a fitted sparse reconstruction, selected native feed-forward channels, and a low-dimensional subspace are different candidate objects, not interchangeable circuits. Appendix~\ref{app:model-background} gives the full model notation and sparse-interface definitions used below.
